\textbf{Phần 1. Trắc nghiệm lựa chọn}
\setcounterpageref{ex}{0}

\begin{ex} %[2H5N27]
Để vẽ một đoạn thẳng dài $7\text{ cm}$, ta cần dùng dụng cụ nào?
\choice
{Com-pa}
{\True Thước kẻ có chia vạch xăng-ti-mét}
{Ê-ke}
{Bút chì màu}
\loigiai{
Dùng thước có chia vạch xăng-ti-mét.\\
\textbf{Đáp án của câu hỏi là B}
}
\end{ex}

\begin{ex} %[2H5N27]
Ghép 2 hình tam giác vuông bằng nhau có thể tạo thành hình nào?
\choice
{\True Hình vuông hoặc hình chữ nhật (hình tứ giác)}
{Hình tròn}
{Hình ngũ giác}
{Đoạn thẳng}
\loigiai{
Ghép 2 tam giác vuông bằng nhau tạo thành hình tứ giác.\\
\textbf{Đáp án của câu hỏi là A}
}
\end{ex}

\begin{ex} %[2H5H27]
Cắt một hình vuông theo đường chéo ta được mấy hình tam giác?
\choice
{$1$ hình}
{\True $2$ hình tam giác}
{$3$ hình tam giác}
{$4$ hình tam giác}
\loigiai{
Cắt theo đường chéo tạo thành 2 hình tam giác.\\
\textbf{Đáp án của câu hỏi là B}
}
\end{ex}

\begin{ex} %[2H5H27]
Muốn vẽ đoạn thẳng $AB$ dài $9\text{ cm}$, ta đặt thước kẻ sao cho vạch số 0 trùng với điểm $A$, điểm $B$ trùng với vạch số
\choice
{$8$}
{$10$}
{\True $9$}
{$1$}
\loigiai{
Điểm $B$ trùng với vạch số 9.\\
\textbf{Đáp án của câu hỏi là C}
}
\end{ex}

\begin{ex} %[2H5H27]
Dùng các que tính dài bằng nhau, để xếp thành 1 hình tam giác cần ít nhất bao nhiêu que tính?
\choice
{$2$ que}
{$4$ que}
{$5$ que}
{\True $3$ que}
\loigiai{
Cần ít nhất 3 que tính.\\
\textbf{Đáp án của câu hỏi là D}
}
\end{ex}

\begin{ex} %[2H5H27]
Dùng các que tính để xếp thành 1 hình vuông cần ít nhất bao nhiêu que tính có độ dài bằng nhau?
\choice
{\True $4$ que}
{$3$ que}
{$5$ que}
{$6$ que}
\loigiai{
Hình vuông có 4 cạnh bằng nhau nên cần 4 que tính.\\
\textbf{Đáp án của câu hỏi là A}
}
\end{ex}

\begin{ex} %[2H5V27]
Để xếp thành 2 hình tam giác rời nhau, cần ít nhất bao nhiêu que tính?
\choice
{$5$ que}
{$7$ que}
{$8$ que}
{\True $6$ que}
\loigiai{
Mỗi tam giác cần 3 que, 2 tam giác rời cần $3 \times 2 = 6$ que tính.\\
\textbf{Đáp án của câu hỏi là D}
}
\end{ex}

\begin{ex} %[2H5V27]
Có 5 que tính bằng nhau, xếp được 2 hình tam giác chung một cạnh. Có tất cả bao nhiêu hình tam giác trong hình vừa xếp?
\choice
{$1$ hình}
{$3$ hình tam giác}
{\True $2$ hình tam giác}
{$4$ hình tam giác}
\loigiai{
Xếp được 2 hình tam giác.\\
\textbf{Đáp án của câu hỏi là C}
}
\end{ex}

\begin{ex} %[2H5V27]
Ghép 4 hình vuông nhỏ cạnh $1\text{ cm}$ lại với nhau thành một hình vuông lớn. Cạnh của hình vuông lớn là
\choice
{$1\text{ cm}$}
{$3\text{ cm}$}
{$4\text{ cm}$}
{\True $2\text{ cm}$}
\loigiai{
Cạnh hình vuông lớn: $1 + 1 = 2\text{ cm}$.\\
\textbf{Đáp án của câu hỏi là D}
}
\end{ex}

\begin{ex} %[2H5V27]
Cắt một tờ giấy hình tứ giác theo một đường thẳng, có thể thu được nhiều nhất mấy hình tam giác?
\choice
{$1$ hình}
{\True $2$ hình tam giác (khi cắt theo đường chéo)}
{$3$ hình}
{$4$ hình}
\loigiai{
Cắt theo đường chéo thu được 2 hình tam giác.\\
\textbf{Đáp án của câu hỏi là B}
}
\end{ex}

\begin{ex} %[2H5V27]
Bạn Nam có 8 que tính bằng nhau. Nam có thể xếp được nhiều nhất bao nhiêu hình vuông rời nhau?
\choice
{\True $2$ hình vuông}
{$1$ hình}
{$3$ hình}
{$4$ hình}
\loigiai{
$8 : 4 = 2$ hình vuông.\\
\textbf{Đáp án của câu hỏi là A}
}
\end{ex}

\begin{ex} %[2H5C27]
Cần ít nhất bao nhiêu que tính để xếp thành 1 hình vuông và 1 hình tam giác chung 1 cạnh?
\choice
{$7$ que}
{$5$ que}
{\True $6$ que ($4 + 3 - 1 = 6$)}
{$8$ que}
\loigiai{
Chung 1 cạnh nên cần: $4 + 2 = 6$ que tính.\\
\textbf{Đáp án của câu hỏi là C}
}
\end{ex}

\textbf{Phần 2. Trắc nghiệm đúng sai}
\setcounterpageref{ex}{0}

\begin{ex} %[2H5H27]
Xét tính đúng sai về xếp ghép hình:
\choiceTF
{\True Ghép hai hình vuông nhỏ bằng nhau có thể tạo thành một hình chữ nhật.}
{\True Một hình vuông có thể cắt thành hai hình chữ nhật bằng nhau.}
{Không thể cắt hình vuông thành các hình tam giác.}
{\True Có thể cắt hình vuông thành 2 hoặc 4 hình tam giác bằng nhau.}
\loigiai{
\begin{itemize}
    \item Ghép 2 hình vuông được hình chữ nhật: Đúng.
    \item Cắt đôi hình vuông được 2 hình chữ nhật: Đúng.
    \item Cắt theo đường chéo ra tam giác được. Ý 3 Sai, Ý 4 Đúng.
\end{itemize}
\begin{itemize}
    \item a) \textbf{Đúng}.
    \item b) \textbf{Đúng}.
    \item c) \textbf{Sai}.
    \item d) \textbf{Đúng}.
\end{itemize}
}
\end{ex}

\begin{ex} %[2H5V27]
Có 7 que tính bằng nhau. Xét tính đúng sai:
\choiceTF
{\True Có thể dùng 7 que tính xếp được 3 hình tam giác chung cạnh.}
{\True Có thể dùng 7 que tính xếp được 1 hình vuông và 1 hình tam giác rời nhau.}
{Không thể xếp được hình tứ giác nào với 7 que tính.}
{\True Có thể xếp được hình tứ giác và thừa que tính.}
\loigiai{
\begin{itemize}
    \item Xếp 3 tam giác chung cạnh cần: $3 + 2 + 2 = 7$ que: Đúng.
    \item 1 hình vuông (4 que) + 1 tam giác (3 que) = 7 que: Đúng.
    \item Dùng 4 que xếp tứ giác vẫn được. Ý 3 Sai, Ý 4 Đúng.
\end{itemize}
\begin{itemize}
    \item a) \textbf{Đúng}.
    \item b) \textbf{Đúng}.
    \item c) \textbf{Sai}.
    \item d) \textbf{Đúng}.
\end{itemize}
}
\end{ex}

\textbf{Phần 3. Trả lời ngắn}
\setcounterpageref{ex}{0}

\begin{ex} %[2H5H27]
Cần bao nhiêu que tính để xếp thành 1 hình tam giác?
\shortans{3}
\loigiai{
\textbf{Đáp số:} $3$\\
3 que tính.
}
\end{ex}

\begin{ex} %[2H5H27]
Cần bao nhiêu que tính bằng nhau để xếp thành 1 hình vuông?
\shortans{4}
\loigiai{
\textbf{Đáp số:} $4$\\
4 que tính.
}
\end{ex}

\begin{ex} %[2H5V27]
Có 9 que tính, xếp được nhiều nhất bao nhiêu hình tam giác rời nhau?
\shortans{3}
\loigiai{
\textbf{Đáp số:} $3$\\
$9 : 3 = 3$ hình.
}
\end{ex}

\begin{ex} %[2H5V27]
Có 12 que tính, xếp được nhiều nhất bao nhiêu hình vuông rời nhau?
\shortans{3}
\loigiai{
\textbf{Đáp số:} $3$\\
$12 : 4 = 3$ hình.
}
\end{ex}

\begin{ex} %[2H5V27]
Cắt một hình chữ nhật bởi 1 nhát kéo thẳng theo đường chéo thì được bao nhiêu hình tam giác?
\shortans{2}
\loigiai{
\textbf{Đáp số:} $2$\\
Được 2 hình tam giác.
}
\end{ex}

\begin{ex} %[2H5C27]
Cần ít nhất bao nhiêu que tính để xếp thành 2 hình vuông có chung 1 cạnh?
\shortans{7}
\loigiai{
\textbf{Đáp số:} $7$\\
$4 + 3 = 7$ que tính.
}
\end{ex}
